\documentclass{article}
\input{../../lib.sty}

\title{\texttt{impl PrivateCandidateMeasure for MaxDivergence}}
\author{Michael Shoemate}
\date{}

\begin{document}
\maketitle

Proves soundness of the \texttt{MaxDivergence} implementation of
\texttt{PrivateCandidateMeasure} in \texttt{mod.rs}.

\section{Hoare Triple}
\subsection*{Precondition}
\texttt{base\_map} is the privacy map of a valid \texttt{MaxDivergence} measurement.

\subsection*{Pseudocode}
\lstinputlisting[language=Python,firstline=2,escapechar=|]{./pseudocode/max_divergence_private_candidate_measure.py}

\subsection*{Postcondition}
\begin{theorem}
\label{postcondition}
The \texttt{validate} method accepts exactly the supported pure-DP combinations implemented in
\texttt{mod.rs}, and \texttt{new\_privacy\_map} returns the privacy map stated in the proof
definition of \texttt{make\_select\_private\_candidate}.
\end{theorem}

\section{Proof}
The validation logic is immediate:
\begin{itemize}
    \item line \ref{line:geometric-guard} rejects thresholded constructions unless the repetition
    family is geometric;
    \item line \ref{line:poisson-guard} rejects all Poisson constructions;
    \item no other combinations are rejected.
\end{itemize}

For privacy accounting, there are two accepted cases.
\begin{itemize}
    \item In the thresholded case, line \ref{line:threshold-factor} sets the factor to $2$.
    By Theorem 3.1 of \cite{liu2018privateselectionprivatecandidates}, the thresholded geometric
    construction is $2\epsilon$-DP when the base mechanism is $\epsilon$-DP.
    \item In the non-thresholded case, line \ref{line:nb-factor} computes an upward-rounded factor
    satisfying \texttt{factor} $\ge 2 + \eta$.
    By Corollary 3 of \cite{papernot2021hyperparametertuningrenyidifferentialprivacy}, the
    negative-binomial best-of-$k$ construction is $(2 + \eta)\epsilon$-DP.
\end{itemize}

Line \ref{line:mul} applies this factor to \texttt{base\_map(d\_in)}, so the returned privacy map
is a conservative upper bound on the implemented pure-DP privacy loss.

\bibliographystyle{plain}
\bibliography{references.bib}

\end{document}
